\begin{ex}
Để vẽ đường thẳng đi qua điểm A và song song với đường thẳng d, ta thường dùng dụng cụ nào kết hợp với thước kẻ? (Thước ê-ke ghi 1, compa ghi 2).
\shortans{1}
\loigiai{
Thước ê-ke (1).
}
\end{ex}

\begin{ex}
Qua một điểm ở ngoài đường thẳng, vẽ được bao nhiêu đường thẳng song song với đường thẳng đó?
\shortans{1}
\loigiai{
Chỉ vẽ được 1 đường thẳng.
}
\end{ex}

\begin{ex}
Để vẽ hình chữ nhật có chiều dài 7 cm và chiều rộng 4 cm, ta cần vẽ các cặp cạnh thế nào với nhau? (Vuông góc và song song ghi 1, chỉ song song ghi 2).
\shortans{1}
\loigiai{
Vuông góc và song song (1).
}
\end{ex}

\begin{ex}
Vẽ hình vuông cạnh 6 cm. Chu vi của hình vuông đó là bao nhiêu xăng-ti-mét?
\shortans{24}
\loigiai{
$6 \times 4 = 24$ cm.
}
\end{ex}

\begin{ex}
Vẽ hình vuông cạnh 6 cm. Diện tích của hình vuông đó là bao nhiêu xăng-ti-mét vuông?
\shortans{36}
\loigiai{
$6 \times 6 = 36\text{ cm}^2$.
}
\end{ex}

\begin{ex}
Vẽ hình chữ nhật có chiều dài 9 cm, chiều rộng 5 cm. Chu vi là bao nhiêu xăng-ti-mét?
\shortans{28}
\loigiai{
$(9 + 5) \times 2 = 28$ cm.
}
\end{ex}

\begin{ex}
Vẽ hình chữ nhật có chiều dài 9 cm, chiều rộng 5 cm. Diện tích là bao nhiêu xăng-ti-mét vuông?
\shortans{45}
\loigiai{
$9 \times 5 = 45\text{ cm}^2$.
}
\end{ex}

\begin{ex}
Hai đường thẳng cùng song song với một đường thẳng thứ ba thì chúng có song song với nhau không? (Có ghi 1, không ghi 0).
\shortans{1}
\loigiai{
Có song song (1).
}
\end{ex}

\begin{ex}
Trên giấy kẻ ô vuông, vẽ hai đường thẳng theo hai dòng kẻ ngang cách nhau 3 ô vuông. Khoảng cách giữa chúng có đổi không? (Có đổi ghi 1, không đổi ghi 0).
\shortans{0}
\loigiai{
Không đổi (0).
}
\end{ex}

\begin{ex}
Vẽ đường thẳng d1 song song với d2 cách nhau 5 cm. Một đường thẳng d3 vuông góc với d1 thì có vuông góc với d2 không? (Có ghi 1, không ghi 0).
\shortans{1}
\loigiai{
Có vuông góc (1).
}
\end{ex}

\begin{ex}
Một hình chữ nhật có chu vi 36 cm, chiều dài gấp đôi chiều rộng. Chiều rộng hình chữ nhật đó là bao nhiêu xăng-ti-mét?
\shortans{6}
\loigiai{
Nửa chu vi = 18 cm. Rộng = $18 : 3 = 6$ cm.
}
\end{ex}

\begin{ex}
Chiều dài hình chữ nhật trong bài toán trên là bao nhiêu xăng-ti-mét?
\shortans{12}
\loigiai{
$6 \times 2 = 12$ cm.
}
\end{ex}

\begin{ex}
Diện tích hình chữ nhật đó là bao nhiêu xăng-ti-mét vuông?
\shortans{72}
\loigiai{
$12 \times 6 = 72\text{ cm}^2$.
}
\end{ex}

\begin{ex}
Một hình vuông có chu vi 32 cm. Diện tích hình vuông đó là bao nhiêu xăng-ti-mét vuông?
\shortans{64}
\loigiai{
Cạnh = 8 cm. Diện tích = $8 \times 8 = 64\text{ cm}^2$.
}
\end{ex}

\begin{ex}
Vẽ một đoạn thẳng AB dài 10 cm, lấy M là trung điểm của AB. Đoạn AM dài bao nhiêu xăng-ti-mét?
\shortans{5}
\loigiai{
$10 : 2 = 5$ cm.
}
\end{ex}

\begin{ex}
Vẽ đường thẳng d qua M vuông góc với AB. Góc tạo bởi d và AB bằng bao nhiêu độ?
\shortans{90}
\loigiai{
$90^\circ$.
}
\end{ex}

\begin{ex}
Khoảng cách giữa hai mép đối diện của chiếc thước kẻ nhựa là 3 cm. Hai mép đó có song song không? (Có ghi 1, không ghi 0).
\shortans{1}
\loigiai{
Có song song (1).
}
\end{ex}

\begin{ex}
Dùng thước vẽ 3 đường thẳng song song cách đều nhau 2 cm. Khoảng cách giữa đường thứ nhất và đường thứ ba là bao nhiêu xăng-ti-mét?
\shortans{4}
\loigiai{
$2 + 2 = 4$ cm.
}
\end{ex}

\begin{ex}
Dùng thước vẽ 4 đường thẳng song song cách đều nhau 2 cm. Khoảng cách giữa đường thứ nhất và đường thứ tư là bao nhiêu xăng-ti-mét?
\shortans{6}
\loigiai{
$2 \times 3 = 6$ cm.
}
\end{ex}

\begin{ex}
Hình thang có 2 cạnh đáy song song dài 6 cm và 10 cm. Chiều cao bằng 4 cm. Diện tích hình thang đó là bao nhiêu xăng-ti-mét vuông?
\shortans{32}
\loigiai{
$(6 + 10) \times 4 : 2 = 32\text{ cm}^2$.
}
\end{ex}

\begin{ex}
Một hình bình hành có độ dài đáy 12 cm, chiều cao 5 cm. Diện tích hình bình hành đó là bao nhiêu xăng-ti-mét vuông?
\shortans{60}
\loigiai{
$12 \times 5 = 60\text{ cm}^2$.
}
\end{ex}

\begin{ex}
Một hình thoi có độ dài hai đường chéo là 8 cm và 6 cm. Diện tích hình thoi đó là bao nhiêu xăng-ti-mét vuông?
\shortans{24}
\loigiai{
$8 \times 6 : 2 = 24\text{ cm}^2$.
}
\end{ex}

\begin{ex}
Hình vuông có cạnh 10 cm. Nối trung điểm 4 cạnh được hình vuông nhỏ. Diện tích hình vuông nhỏ là bao nhiêu xăng-ti-mét vuông?
\shortans{50}
\loigiai{
Diện tích = $100 : 2 = 50\text{ cm}^2$.
}
\end{ex}

\begin{ex}
Hai đường thẳng song song có cắt nhau ở vô cực không? (Không ghi 0, có ghi 1).
\shortans{0}
\loigiai{
Không cắt nhau (0).
}
\end{ex}

\begin{ex}
Số góc vuông trong một hình vuông là bao nhiêu?
\shortans{4}
\loigiai{
Có 4 góc vuông.
}
\end{ex}
